\section{Sviluppi Futuri}
Per ragioni progettuali, di tempo e date le forti limitazioni, sono rimasti alcuni punti che potrebbero essere modificati con lo scopo di migliorare la stabilità e il rendimento del progetto.

\begin{itemize}
\item   Reset asincrono: Data la specifica, per la quale non si potevano inserire componenti all’esterno delle Butterfly,  abbiamo una connessione che risulta diretta con l’esterno e quindi condizionata dal rumore. La presenza di glitch rischierebbe quindi di resettare la macchina perdendo il lavoro in esecuzione. Questo problema può essere risolto con il campionamento di suddetto ingresso o anche solo con un controllo della durata del reset con lo scopo di evitare che brevissimi glitch influenzino la macchina a stati (una porta AND fra il segnale di reset e il segnale di reset ritardato tramite 2$\cdot$N inverter in cascata, il cui numero è da stabilire in base alla tecnologia).
\item   Gestione dei W : Siamo consapevoli che la soluzione meno costosa e più efficace è sicuramente quella di porre i pin di ingresso per ogni singola Butterfly ai suoi $W_r$ e $W_i$ dedicati, ma questo può limitare la flessibilità del progetto e, in caso di errore nella realizzazione dei valori inchiodati, potrebbe rendere inutilizzabile il circuito. Perciò, se servisse maggiore flessibilità nell'invio dei W, si è quindi pensato a un sistema basato sul passaggio dei W attraverso un Bus e un enable codificato per la selezione dei riceventi. Questa soluzione però è soddisfacente fino al caso con 32 campioni, dato che con 64 campioni (per cui servono 6 colonne da 32 Butterfly) si avrebbero 32 W da passare per un totale di 64 valori che vanno consegnati entro lo start dell’ultima colonna che avviene al 55-esimo colpo (T\_ Butterfly * (N\_colonne -1)=11*5=55).
\item Uscite: per limitare il rumore e i consumi prodotti dalle uscite delle singole Butterfly si potrebbe introdurre una porta AND per ogni uscita, che fa passare l’uscita solo quando si sta trasmettendo un segnale DONE e al colpo ad esso successivo (salvando questa informazione in un FF).
\end{itemize}
